\documentclass[11pt]{article}

\usepackage[T1]{fontenc}
\usepackage[utf8]{inputenc}
\usepackage{lmodern}
\usepackage[margin=1.04in]{geometry}
\usepackage{amsmath,amssymb,mathtools,amsthm}
\usepackage{aliascnt}
\usepackage{mathrsfs}
\usepackage{microtype}
\usepackage{enumitem}
\usepackage{xcolor}
\usepackage{booktabs}
\usepackage{url}
\usepackage{hyperref}
\usepackage[capitalize,noabbrev]{cleveref}

\hypersetup{
  colorlinks=true,
  linkcolor=blue!45!black,
  citecolor=green!35!black,
  urlcolor=blue!55!black,
  pdftitle={Multicolumn Exponential Periods: Conditional Comparison and Geometric Realization},
  pdfauthor={Christopher D. Long; Antoine-Auguste Le Blanc},
  pdfsubject={Candidate multicolumn affine-orbit and actual-realization program},
  pdfkeywords={exponential periods, candidate draft, comparison, open problems}
}

\numberwithin{equation}{section}
\setlength{\emergencystretch}{3em}

\newtheorem{theorem}{Theorem}[section]

\newaliascnt{proposition}{theorem}
\newtheorem{proposition}[proposition]{Proposition}
\aliascntresetthe{proposition}

\newaliascnt{lemma}{theorem}
\newtheorem{lemma}[lemma]{Lemma}
\aliascntresetthe{lemma}

\newaliascnt{corollary}{theorem}
\newtheorem{corollary}[corollary]{Corollary}
\aliascntresetthe{corollary}

\theoremstyle{definition}
\newaliascnt{definition}{theorem}
\newtheorem{definition}[definition]{Definition}
\aliascntresetthe{definition}

\newaliascnt{example}{theorem}
\newtheorem{example}[example]{Example}
\aliascntresetthe{example}

\theoremstyle{remark}
\newaliascnt{remark}{theorem}
\newtheorem{remark}[remark]{Remark}
\aliascntresetthe{remark}

\newcommand{\C}{\mathbb C}
\newcommand{\Q}{\mathbb Q}
\newcommand{\Qbar}{\overline{\mathbb Q}}
\newcommand{\R}{\mathbb R}
\newcommand{\Z}{\mathbb Z}
\newcommand{\N}{\mathbb N}
\newcommand{\Aone}{\mathbb A^1}
\newcommand{\Atwo}{\mathbb A^2}
\newcommand{\Gm}{\mathbb G_m}
\newcommand{\Spec}{\operatorname{Spec}}
\newcommand{\Span}{\operatorname{span}}
\newcommand{\trdeg}{\operatorname{trdeg}}
\newcommand{\per}{\operatorname{per}}
\newcommand{\ord}{\operatorname{ord}}
\newcommand{\End}{\operatorname{End}}
\newcommand{\Hom}{\operatorname{Hom}}
\newcommand{\Ext}{\operatorname{Ext}}
\newcommand{\Gal}{\operatorname{Gal}}
\newcommand{\Lie}{\operatorname{Lie}}
\newcommand{\Frac}{\operatorname{Frac}}
\newcommand{\Sym}{\operatorname{Sym}}
\newcommand{\SL}{\operatorname{SL}}
\newcommand{\GL}{\operatorname{GL}}
\newcommand{\Pic}{\operatorname{Pic}}
\newcommand{\Div}{\operatorname{Div}}
\newcommand{\Res}{\operatorname{Res}}
\newcommand{\Tr}{\operatorname{Tr}}
\newcommand{\Crit}{\operatorname{Crit}}
\newcommand{\supp}{\operatorname{supp}}
\newcommand{\rank}{\operatorname{rank}}
\newcommand{\Ecal}{\mathcal E}
\newcommand{\Fcal}{\mathcal F}
\newcommand{\Hcal}{\mathcal H}
\newcommand{\Lcal}{\mathcal L}
\newcommand{\Mcal}{\mathcal M}
\newcommand{\Ocal}{\mathcal O}
\newcommand{\Pcal}{\mathcal P}
\newcommand{\Rcal}{\mathcal R}
\newcommand{\Vcal}{\mathcal V}
\newcommand{\dd}{\,d}
\newcommand{\KZ}{\mathrm{KZ}}

\title{Multicolumn Exponential Periods: Conditional Comparison and Geometric Realization}
\author{Christopher D. Long \and Antoine-Auguste Le Blanc\thanks{Le Blanc cryptographic author identifier: \texttt{b6048125\allowbreak{}30b3535b\allowbreak{}23d6dbb0\allowbreak{}f5abe32d\allowbreak{}1ed1a9de\allowbreak{}35fb4f7d\allowbreak{}04d62cff\allowbreak{}08ea19c8}}}
\date{October 2026}


\newcommand{\Hq}{H_q^2}
\newcommand{\M}{\mathscr M}
\newcommand{\Ncal}{\mathcal N}
\newcommand{\Tcal}{\mathcal T}
\newcommand{\A}{\mathbb A}

\newcommand{\PGL}{\operatorname{PGL}}
\newcommand{\wt}{\operatorname{wt}}
\newcommand{\Bcal}{\mathcal B}

\newtheorem{conjecture}[theorem]{Conjecture}
\begin{document}
\maketitle

\begin{center}
\small
\texttt{galizur@gmail.com}
\end{center}

\begin{abstract}
This candidate draft isolates the algebraic orbit calculations useful for
multicolumn exponential-period comparison and distinguishes them from their
unproved realization by actual collections of chains. We prove the
rectangular and square special-linear orbit statements, explain joint
determinant-character relations, and give an explicit conditional descent
criterion. The missing geometric and arithmetic hypotheses are listed
rather than assigned to a formal period category by definition. No
unrestricted multichain numerical comparison theorem is asserted.
\end{abstract}
\begin{center}
\textbf{Candidate research draft; not a numbered paper in the series.}
\end{center}
\tableofcontents

\section{Why a single column is different}
Paper V \cite{PaperV} studies a chosen actual vector $U$ satisfying
$U'=AU+b$ over a fixed boundary field. An affine differential Galois group
can have a full translation kernel, which gives a dense affine orbit. If
the kernel is zero, a reductive cocycle splits, giving
\[
 U=a+Yc,
\]
where $Y$ is a fundamental matrix and $c$ is a constant vector. For a
homogeneous group containing $\SL_n$, one nonzero vector still has dense
orbit when $n\ge2$.

Several chains produce several columns. Even when every individual column
passes the one-column test, their combined orbit can satisfy determinant
relations. Thus a product of individually nonzero scalar periods, as used
in Paper VI \cite{PaperVI}, is not a multicolumn independence theorem.

\section{Explicit special-linear orbits}
Let $C$ be an algebraically closed field of characteristic zero, let
$n\ge2$, and let $\SL_n(C)$ act on $n\times r$ matrices by left
multiplication.

\begin{proposition}[Rectangular columns]\label{prop:rectangular}
If $1\le r<n$ and a matrix $Z$ has column rank $r$, its $\SL_n$ orbit
is the set of all rank-$r$ matrices. This orbit is Zariski dense in
$\A^{nr}$ and has zero polynomial vanishing ideal.
\end{proposition}
\begin{proof}
Given two ordered independent $r$-tuples, extend them to bases of $C^n$.
There is a change of basis carrying one tuple to the other. Because an
unused basis vector exists, rescale its image to make the determinant
one without changing the prescribed first $r$ columns. The full-rank
locus is a nonempty open subset of affine space, hence dense.
\end{proof}

\begin{proposition}[Square columns]\label{prop:square}
If $r=n$ and $Z$ is invertible, its orbit is
\[
 \{X:\det X=\det Z\}.
\]
The polynomial ideal of the orbit is generated by
$\det X-\det Z$.
\end{proposition}
\begin{proof}
For any $X$ with the same nonzero determinant, $XZ^{-1}\in\SL_n$.
The fiber is a translate of $\SL_n$, hence irreducible and reduced.
For completeness, the differential of the determinant at an invertible
matrix is nonzero, so the fiber is smooth. Its defining principal ideal
is therefore radical, and the displayed set equality gives the ideal.
\end{proof}
For example, two independent columns for $\SL_2$ have the determinant
invariant $x_{11}x_{22}-x_{12}x_{21}$, although either column alone has
dense orbit in $\A^2$. This example is an algebraic group calculation,
not a realized pair of exponential chains.

\begin{remark}[Dependent and excess columns]
If the actual constant column matrix has rank less than $r$, its linear
column relations must be imposed first. If $r>n$, maximal minors and their
relations enter; rectangular density cannot be applied without a rank
reduction. These cases require the correct actual constant matrix, not a
chosen formal substitute.
\end{remark}

\section{Several blocks and joint determinant characters}
Suppose distinct homogeneous blocks have derived groups containing a
product of special-linear groups. For a full-rank square block $s$, put
$\delta_s=\det Z_s$. Its residual transformation is governed by a
character $\chi_s$ of a common algebraic group $G$.

\begin{proposition}[Joint character equations]\label{prop:lattice}
For the morphism
\[
 \chi:G\longrightarrow(\Gm)^b,\qquad
 g\longmapsto(\chi_1(g),\ldots,\chi_b(g)),
\]
let $T=\chi(G)$ and
$\Lambda=\{a\in\Z^b: \chi_1^{a_1}\cdots\chi_b^{a_b}=1\}$.
The Laurent ideal of the orbit $\delta T$ is
\[
 (X^a-\delta^a:a\in\Lambda)
 \subset C[X_1^{\pm1},\ldots,X_b^{\pm1}].
\]
\end{proposition}
\begin{proof}
A closed subgroup of a torus is defined by the characters trivial on it.
Equivalently its coordinate algebra is the group algebra
$C[\Z^b/\Lambda]$, and the quotient map on Laurent monomials has kernel
generated by $X^a-1$, $a\in\Lambda$. Translation by $\delta$ gives the
formula. This also covers a disconnected image: in characteristic zero
the finite diagonalizable part is reduced.
\end{proof}
Independent determinant equations for the separate blocks need not
suffice. For example the image $\{(t,t):t\in C^\times\}$ has neither
coordinate fixed but has the joint equation
$X_1\delta_2-X_2\delta_1=0$.

The archived pre-single-chain manuscript \cite{Archive} develops a more
general affine-orbit calculation, including active translation columns,
reductive splitting and local determinant-character models. The present
note does not reclassify arbitrary isotypic affine extensions. In
particular the hypothesis that a derived group is a product of
independent special-linear factors must be proved for the actual blocks.

\section{A conditional actual specialization criterion}
The following algebraic formulation isolates what is required after an
actual functional comparison has been made. A discrete valuation ring
is denoted by $R$, with uniformizer $\pi$, fraction field $K$, and residue
field $\kappa$.

\begin{proposition}[Local contraction and specialization]\label{prop:local}
Let $I\subset R[X_1,\ldots,X_N]$ and suppose
$A=R[X]/I$ is $R$-flat. Then
\[
 I=(I K[X])\cap R[X].
\]
Suppose further that an actual functional period presentation has generic
kernel $I K[X]$, and that every numerical relation at the chosen parameter
lifts to a functional relation with coefficients in $R$ and the prescribed
residue row or polynomial. Then its numerical kernel is exactly the
specialization of $I$, provided every specialized generator of $I$ is an
actual numerical relation.
\end{proposition}
\begin{proof}
Flatness implies torsion-freeness over the discrete valuation ring.
If the class of $P\in R[X]$ vanishes after tensoring with $K$, it is
annihilated by a power of $\pi$, hence already zero. This proves
contraction. A lifted numerical relation is in the generic kernel and
therefore in $I$ by contraction; reducing modulo $\pi$ proves one kernel
inclusion. The stipulated validity of specialized generators proves the
other inclusion.
\end{proof}
The assertion is conditional on the actual presentation and the lifting
hypothesis. It does not prove those hypotheses. In a period application,
$R$ must be chosen from actual boundary coefficients regular at the
specified parameter; a formal localization at unspecified denominators
is insufficient. Beukers' lifting theorem \cite{Beukers} applies to
closed $E$-function systems at nonsingular nonzero algebraic arguments,
not automatically to an arbitrary formal coefficient ring.

\section{What a genuine multichain theorem must verify}
A publishable realization should exhibit specific chains and prove the
following items in dependency order.
\begin{enumerate}[leftmargin=2em]
\item \emph{Actual columns and common base.} Derive the common inhomogeneous
system by literal polynomial Stokes certificates, and identify all actual
linear relations among columns. Passing to a cover cannot silently be
assumed to preserve irreducibility or independence of blocks.
\item \emph{Homogeneous and affine groups.} Verify the product derived group,
its surviving boundary-field action, and the full translation subspace.
Isomorphic blocks require an isotypic analysis, not separate copies of an
irreducible representation by fiat.
\item \emph{Determinant realization.} Construct determinant-line relations
with the actual period normalizations. An invariant of an abstract
fundamental matrix does not by itself identify a formal Stokes relation
between the prescribed chains.
\item \emph{Joint characters and local units.} Compute the entire character
lattice and prove that determinant coefficients and affine shifts have
the required regularity and nonvanishing at specialization.
\item \emph{Numerical lifting.} Place the actual coordinates in a closed
$E$-function system and verify the ordinary-point hypotheses and local
contraction used in \cref{prop:local}.
\end{enumerate}
The earlier unrestricted endpoint-complete claim in the development
history was withdrawn. Neither this candidate nor the new scalar theorem
in Paper V reinstates it.

\section{A concrete first target and stopping rule}
The smallest informative example is a rank-two interior system with two
geometrically different compact chains sharing a specified boundary
system. A successful example should first determine whether the affine
translation kernel makes both columns free or whether a normalized
determinant survives. In the latter case the determinant must be computed
as an actual boundary quantity and checked at every claimed algebraic
scale. A symbolic $2\times2$ determinant calculation is only the orbit
part of this task.

The candidate is deliberately not assigned a series number. Its proved
statements are the elementary orbit and local-contraction propositions
above. The missing theorem is an actual multichain realization satisfying
all five requirements, with at least one fully checked example. The
archived proofs and the current V should be cited at their exact scope,
not used as evidence that the realization has already been completed.

\begin{thebibliography}{9}
\raggedright
\bibitem{PaperV} C. D. Long and A.-A. Le Blanc, \emph{Relative Exponential
Periods on Polynomial Chains: Algebraic Independence over Fixed Boundary
Fields}, Paper V, September 2026, revised October 2026.
\bibitem{PaperVI} C. D. Long and A.-A. Le Blanc, \emph{Exponential Periods
and Factorial Moments}, Paper VI, October 2026, working manuscript.
\bibitem{Archive} C. D. Long and A.-A. Le Blanc, \emph{Relative Exponential
Periods on Polynomial Flags: Polynomial Boundary Systems and an Elliptic
Comparison Theorem}, pre-single-chain manuscript, archived 6 October 2026,
appendix on affine orbit ideals and local flat models.
\url{https://github.com/long-mathematics/exponential-integral-theorem/tree/main/notes/archives/2026-10-06-paper-v-before-single-chain}.
\bibitem{Beukers} F. Beukers, \emph{A refined version of the
Siegel--Shidlovskii theorem}, Ann. of Math. (2) \textbf{163} (2006), 369--379.
\url{https://annals.math.princeton.edu/wp-content/uploads/annals-v163-n1-p08.pdf}.
\end{thebibliography}
\end{document}
